\documentclass[../DoAn.tex]{subfiles}
\begin{document}
\section{Kết luận}
RoomHub hiện thực bài toán đặt phòng học, phòng họp trên Android. Hệ thống gồm ứng dụng di động, dịch vụ Java/Spring Boot, REST API và cơ sở dữ liệu H2 được truy cập qua JPA. Cấu hình hiện tại chọn chế độ dùng máy chủ; chế độ lưu cục bộ vẫn phục vụ trình diễn trên một thiết bị. Bản phát hành Android đã được kiểm tra chạy độc lập ở chế độ cục bộ.

Các kết quả kiểm thử cần được hiểu theo đúng phạm vi. Kiểm thử ứng dụng bao phủ nhiều quy tắc ở chế độ cục bộ; dịch vụ Java có kiểm thử tích hợp. Luồng trao đổi dữ liệu giữa ứng dụng và máy chủ đã được lập trình, nhưng vận hành thực tế phụ thuộc vào máy chủ, cơ sở dữ liệu và địa chỉ API có thể truy cập. Đường kết nối công khai hiện chỉ dùng cho thử nghiệm. Dữ liệu cục bộ không tự hợp nhất với dữ liệu máy chủ khi mất mạng; khả năng chia sẻ dữ liệu giữa nhiều thiết bị cần được xác nhận đầu cuối. Với khóa số, việc gửi bản tin OneIoT chưa chứng minh khóa thật đã nhận hoặc áp dụng mã.

\section{Hướng phát triển}
Nhóm cần kiểm tra đầu cuối chế độ máy chủ với nhiều thiết bị, bổ sung kiểm thử đồng thời cho các trường hợp đặt cùng phòng và làm rõ cơ chế xử lý xung đột khi mạng gián đoạn. Khi triển khai lâu dài, cần rà lại lược đồ dữ liệu, bảo vệ mật khẩu tài khoản và mã khóa, quản lý phiên truy cập, và kiểm tra bảo mật của kết nối OneIoT. Một bước quan trọng khác là nhận phản hồi từ khóa thật để xác nhận trạng thái tạo mã, check-in và check-out thay vì chỉ dựa vào trạng thái gửi bản tin.

Về nghiệp vụ, có thể mở rộng quản lý nhiều khóa, theo dõi chi tiết thiết bị trong phòng, ghi nhật ký thao tác và xây dựng báo cáo sử dụng. Mỗi chức năng mới cần được kiểm chứng ở cả chế độ cục bộ và chế độ dùng máy chủ trước khi coi là hoàn thiện.
\end{document}
